
\section{Overview}
The supplementary materials is structured as follows:
\begin{itemize}
    \item \textbf{\autoref{sec:appendix-limitations}:~\nameref{sec:appendix-limitations}}.
    \item \textbf{\autoref{sec:appendix_details_discoveryworld}:~\nameref{sec:appendix_details_discoveryworld}}. This section covers the \textit{\nameref{sec:appendix_action_space}} and common \textit{\nameref{sec:appendix-object-properties}} of \env, the recommended practices for different \textit{\nameref{sec:train_eval_splits}}, more details on the \textit{\nameref{sec:appendix-discovery-themes}} and the \textit{\nameref{sec:appendix_unit_tests}}.
    
    \item \textbf{\autoref{sec:appendix-additional-model-details}:~\nameref{sec:appendix-additional-model-details}}.

    \item \textbf{\autoref{sec:appendix-automatic-knowledge-evaluation}:~\nameref{sec:appendix-automatic-knowledge-evaluation}}.

    \item \textbf{\autoref{sec:appendix_human}:~\nameref{sec:appendix_human}}.
\end{itemize}

\clearpage

\section{Limitations and Broader Impacts}
\label{sec:appendix-limitations}

\paragraph{Simulation Fidelity:} \env is inherently a low-fidelity representation of the physical world, with an abstracted action space.  As such, agents that perform well on \env may not necessarily perform well at making discoveries in the real world, given the much larger search and action space.  That being said, while the simulated environment is low-fidelity compared to the physical world, the steps of discovery that are simulated (from ideation, hypothesis formation, experiment design and execution, data analysis, and forming conclusions) are common steps in the scientific method regardless of whether those actions are taking place in the real world or a virtual environment. 

\paragraph{Agent Cost:} The agent models that we run are currently quite costly, ranging from approximately USD\$3k-\$10k to run for the complete set of 120 tasks in \env.  This cost is due in large part to (1) the long (1000+ steps) runtimes, each step requiring at least one LLM call, (2) the large number of individual tasks to evaluate, and (3) the use of performant but costly \textsc{API-based} models that charge per token.  We believe that developing inexpensive models that allow for rapid iteration is a clear near-term goal for developing further discovery agents. 

\paragraph{Societal Benefits and Risks:} Automated scientific discovery has the potential for broadly positive societal impact by accelerating the pace of scientific progress, potentially decreasing the time it takes for novel discoveries in medicine, chemistry, technology, and other areas of broad impact.  If discovery systems are used by individuals or groups with prosocial intentions, there is the potential for broad societal impact.  Conversely, if individuals or groups with negative intentions choose to attempt to accelerate discoveries that may be harmful, there is the potential for those individuals or groups to cause harm using automated scientific discovery models.

\section{Additional Details on \env}
\label{sec:appendix_details_discoveryworld}

To encourage the community to work on general AI discovery agents, we have open-sourced \env under Apache-2.0 license on GitHub at \href{https://github.com/allenai/discoveryworld}{github.com/allenai/discoveryworld}.

\subsection{Action Space}
\label{sec:appendix_action_space}

The action space for \env is shown in Table~\ref{tab:action_space}.

\begin{table}[ht]
\centering
\footnotesize
\caption{\footnotesize The action space for DiscoveryWorld, which includes 14 main actions.  Actions can take zero (e.g. \textsc{wait}), one (e.g. \textsc{take thermometer}), or two (e.g. \textsc{use spectrometer on soil}) arguments. Two handicap actions ($*$) are available to agents to assist with poor spatial/navigation abilities.}
\label{tab:action_space}
\begin{tabular}{p{3.1cm}p{3cm}p{2.7cm}p{3.2cm}}
\toprule
\textbf{Action} &   \textbf{Description}                        &   \textbf{Action} &   \textbf{Description}    \\
\midrule
\green{Move} \blue{DIR} & Move North                            &   \green{Use} \blue{OBJ} \green{[on} \blue{OBJ}\green{]} & Use spectrometer on soil        \\
\green{Take} \blue{OBJ} & Take thermometer                      &   \green{Eat} \blue{OBJ} & Eat mushroom        \\
\green{Drop} \blue{OBJ} & Drop seed                             &   \green{Read} \blue{OBJ} & Read rocketry book         \\
\green{Put/Give} \blue{OBJ} \green{to} \blue{OBJ}    & Put sample in jar  &   \green{Wait} & Do nothing        \\
\green{Open/Close} \blue{OBJ} & Open door                       &   \green{Feed} & View DiscoveryFeed        \\
\green{De/Activate} \blue{OBJ} & Activate pendulum              &   \green{Teleport} \blue{LOC} $^{*}$ & Teleport to Science Lab        \\
\green{Talk} \blue{OBJ} & Talk to Colonist                      &   \green{Teleport} \blue{OBJ} $^{*}$ & Teleport to Microscope         \\
\bottomrule        
\end{tabular}
\end{table}

\subsection{Experimental Settings}
\label{sec:train_eval_splits}

As both a development environment and benchmarking tool, we provide the following recommended configurations for evaluating on \env in common experimental settings:

\paragraph{Zero-shot:} For zero-shot evaluation, an agent should not have prior exposure to \env, and can be evaluated on all 120 tasks.  Each task should be evaluated \textit{independently}, without any carry-over knowledge from one task to another.  \textit{All of the baseline agent models provided in this work are evaluated zero-shot.}

\paragraph{Single-Task Learning:} For agents that require training data (e.g., RL agents, few-shot examples, etc.) of highly similar tasks (i.e. within-task designs), following the \textsc{TextWorld}~\citep{Ct2018TextWorldAL}, \textsc{AlfWorld}~\citep{Shridhar2020ALFWorldAT}, and \textsc{ScienceWorld}~\citep{wang2022scienceworld} conventions, we recommend the following \textit{within-theme} design: training on parametric seeds 0 and 1, developing on seed 2, and testing on seeds 3 and 4.  In addition, all tasks in the \textit{Unit Tests} are available for training.

\paragraph{Multi-Task Learning:} For agents that require training data (e.g., RL agents, few-shot examples, etc.), following the \textsc{Jericho} benchmark~\citep{Hausknecht2019InteractiveFG} for evaluating text game suites, we recommend training and evaluating \textit{across themes}.  Tasks 1 to 6 (i.e., theme \textit{Proteomics} and \textit{Chemistry}) can be used for training, tasks 7 through 12 (\textit{Archaeology} and \textit{Reactor Lab}) for development, and the remaining tasks 13 through 24 (\textit{Plant Nutrients, Space Sick, Rocket Science}, and \textit{Translation}) for testing.  In addition, all tasks in the \textit{Unit Tests} are available for training.

\paragraph{Curriculum Learning:} For agents that require training data of growing complexity, irrespective of the experimental setting, we recommend to simply follow the natural progression of the difficulty levels for each theme. In addition, all tasks in the \textit{Unit Tests} are available for training.



\subsection{Discovery Themes}
\label{sec:appendix-discovery-themes}

The following are descriptions of the 8 main discovery themes, and how these themes vary across each of the 3 levels of difficulty.  \textit{Note: This section contains spoilers!}

\paragraph{Proteomics} This theme examines an agent's ability to perform analyses that involve clustering and identifying outliers.  The agent is told that a number of animal species exist in a remote habitat, with one hypothesized to have migrated from an isolated island in the distant past -- with the agent required to measure the abundance of different proteins (using a \textit{protein meter}) to determine which species is the outlier.  In \textit{normal} difficulty the clustering problem is two dimensional (that is, values of two proteins are provided for each measured animal), while in \textit{challenge} difficulty, the clustering problem is three dimensional. Protein concentrations are provided as values between zero and one. The protein concentration data is manufactured such that all animals protein values exist on the surface of a sphere centered on a randomly chosen point, with inliers having a small radius from the center (e.g. $0.1$), and outliers having a much farther radius (e.g. $0.4$).  Inliers and outliers are clearly visible on a plot by inspection.  The \textit{easy} difficulty level reduces the number of animals from 5 to 3, and has the entire task take place in a small lab with unmoving animals, rather than a large map with moving animals.

\paragraph{Chemistry} The chemistry theme evaluates an agents ability to explore a large solution space in the presense of a hill-climbing signal.  Here, the agent is told they must manufacture a chemical (a rust removal agent), from some combination of unique chemicals available to them.  The agent must enumerate the possible combinations of chemicals, prepare these solutions, and test them on a specific rusted object.  After being immersed in a solution, the rusted object will decrease from \textit{heavily rusted} to \textit{moderately} or \textit{lightly rusted} based on the cosine similarity between the chemical concentration of the solution that it is immersed in, and the correct solution, providing a hill-climbing signal for agents to use to more quickly explore the space of possible solutions.  The \textit{normal} difficulty has three possible chemical dispensers, the \textit{challenge} has four chemical dispensers, and both require agents to find specific concentrations of chemicals (e.g. \textit{1 part Chemical A, 2 parts Chemical B}).  The \textit{easy} difficulty has four chemical dispensers, but requires no mixing -- one of the four chemicals is a direct solution.

\paragraph{Archaeology} The archaeology theme involves measuring (and validating) artifacts using radioisotope dating.  In the \textit{challenge} difficulty, the agent is presented with a dig site that contains six artifacts -- three known, three unknown -- as well as a radioisotope meter that measures 4 different radioisotopes.  The agent is prompted that it is unknown if radioisotope dating works on \textit{Planet X}, and if it does, which of the 4 radioisotopes will be useful for dating.  The agent is tasked with finding the oldest unknown artifact, and placing a red flag beside its dig site.  The known artifacts include one from the stone age (a stone hammer), one from the bronze age (a bronze chisel), and one from the iron age (iron tongs).  The agent must know that in radioisotope dating, older artifacts have lower concentrations of a given radioisotope due to radioactive decay.  The agent must also make the critical insight that stone age artifacts are older than bronze age artifacts, and iron age artifacts are younger than bronze age artifacts.  The data are manufactured such that one radioisotope follows this pattern, while others don't (with a correlation of $R^{2} < 0.1$ between age and radioisotope value for all incorrect channels).  In the \textit{normal} difficulty, the instrument directly supplies age, and evaluates instrument use rather than instrument validation.  The \textit{easy} difficulty is similar to \textit{normal} but in a small lab environment.

\paragraph{Reactor Lab} The reactor lab theme requires agents to discover and use mathematical relationships between measured quantities using regression.  The reactor lab contains a number of \textit{quantum crystals}, each with a specific resonance frequency.  The agent must place all crystals into specific slots, and tune each slot to that crystal's frequency, for the reactor to activate.  The frequency depends upon one of five measured properties -- density, temperature, radioactivity, crystal size, or crystal spectra -- which can be measured using instruments in the adjoining science lab.  The agent is provided with some number of crystals with known frequencies, and two crystals with unknown frequencies.  They must measure each crystal with each of the instruments, determine which physical property is related to the frequency (as well as how), and use this equation to calculate the frequencies of the unknown crystals.  In the \textit{normal} setting, 2 known crystals are provided, and the relationship can be found with linear regression ($y=mx+b$).  In the \textit{challenge} setting, 3 known crystals are provided, and the relationship is quadratic ($y=a^2x+bx+c$).  In the \textit{easy} setting, only a single (correct) instrument is provided, there is only one unknown crystal, and the relationship requires only inferring the slope ($y=mx$).

\paragraph{Plant Nutrients} The plant nutrient theme requires agents to infer systems of rules based on observing both positive and negative examples.  The agent finds itself at a botanical research center, and given the task of identifying the (unusual) nutrients that plants on \planetX prefer.  The research center includes a ``pilot field`` where 12 seeds have been planted, some of which sprouted correctly, as well as 3 ``test fields``, each controlled by a soil computer that allows manipulating the nutrients in its soil.  The agent must infer the necessary levels of nutrients in the soil from the pilot field, then successfully grow at least two new plants by setting the nutrient levels in a test field to be appropriate, and planting (and growing) seeds in that field.  All settings include five different nutrients.  In the \textit{normal} setting, rules are simple presence-at-value rules (i.e. plants require a specific nutrient, at a specific value of either \textit{low, medium,} or \textit{high} to grow), and these can be inferred from positive examples alone.  In the \textit{challenge} setting, the rules involve logical relations (e.g. \textsc{XOR}, \textsc{AND}, \textsc{OR}, \textsc{NOT}) between nutrients, and more examples are provided.  The \textit{easy} setting resembles \textit{normal} except that nutrients are binary (\textit{present/absent}) rather than varying in concentration.

\paragraph{Space Sick} The space sick theme requires agents to use open-ended discovery skills to investigate the cause of an illness.  The agents are tasked with identifying why some colonists are becoming mildly ill when eating local food, and correcting this.  The agent must discover that some food has been contaminated by mold, which can be directly observed with some instruments, and indirectly observed with others (e.g. through elevated spectrometer readings), or indirectly inferred (e.g. cooking food causes it to become safe).  Distractors (such as some food being slightly radioactive) lead the agent down paths that do not solve the task.  In the \textit{normal} difficulty, the mold is directly observable using instrumentation.  In the \textit{challenge} difficulty, the agent must discover novel detection instrumentation.  In the \textit{easy} setting, the agent is placed in a lab with 3 samples of food and 4 instruments, and must identify which food sample is contaminated (which is detectable by only a single instrument). 

\paragraph{It's (not) Rocket Science!} The rocket science theme requires agents to measure quantities in an environment, then substitute these quantities into known equations to complete a task.  Here, the agent is tasked with sending a rocket into a specific orbital height around \planetX. In the \textit{challenge} version of this task, the agent needs to enter to appropriate orbital velocity, as well as the type (and quantity) of propellant to be used. The agent is provided with a rocketry book containing all the formulas needed to complete the task. However, before it can calculate the required quantities, it needs to infer out how to measure unknown values such as \planetX's gravity and radius, as well as properties related to the three available types of propellent such as density, mass flow rate, and thrust when consumed by the rocket engine.  This theme involves recalling known experiments performed on Earth and transposing them to \planetX (e.g., Eratosthenes' classical experiment to measure Earth's radius\footnote{\href{https://en.wikipedia.org/wiki/Eratosthenes\#Measurement_of_Earth's_circumference}{https://en.wikipedia.org/wiki/Eratosthenes\#Measurement\_of\_Earth's\_circumference}}).  The \textit{normal} difficulty simplifies the task to requiring only orbital velocity (not propellant), while the \textit{easy} difficulty tells the agent they are on a planet with a similar mass and radius to a known celestial body (e.g. Earth, Mars, Venus), eliminating the requirement to measure these values of \planetX.

\paragraph{Lost in Translation} The translation theme requires agents to infer the meanings of progressively more complicated utterances given access to an environment containing Rosetta-stone-style information.  Here, the agent must translate an unknown language used by the native inhabitants of \planetX. The agent first talks to one of the inhabitants, who speaks an utterance.  From this, the agent must explore the village to gather clues about what each word means (e.g., visiting shops will help identifying the name of the items being sold there). In the \textit{challenge} difficulty, instructions are composed of a verb, an amount, a color, and an item name (e.g., bring 3 red flowers). Bringing the correct number of the correct objects back to the inhabitant will complete the task.  In the \textit{normal} difficulty, the utterance requires translating only a single specific object.  In the \textit{easy} setting, the task is reduced to a small lab setting, where the agent must identify which item to bring to another agent based on its translation (accessible by a sign near the object).



\subsection{Unit Tests}
\label{sec:appendix_unit_tests}

\begin{figure}[h!]
  \centering
  \includegraphics[scale=0.5]{figures/scenario-unit-tests.pdf}
  \caption{\footnotesize Example instances of the 10 Unit Test themes.}
  \label{fig:unit_test_screenshots}
\end{figure}

\begin{table}[t!]
\centering
\scriptsize
\caption{\footnotesize High-level descriptions of the 10 \textit{unit test themes} in \env.}  Specific unit tests tasks are parametrically generated from a given theme.
\label{tab:unit-task-descriptions}
\begin{tabular}{p{0.3cm}p{3.7cm}p{8cm}}
\toprule
\#  &   \textbf{Theme} & \textbf{Description} \\
\midrule
\rowcolor[HTML]{F3F3F3} 
25 & Multi-turn dialog with agent & The user agent must talk to an NPC agent, and respond correctly to its requests.  The NPC agent asks the user agent to select the dialog item that says (for example) `kiwi`, and the user agent must select this dialog option in the dialog tree amongst 8 distractors.  The task completes when this process happens successfully 3 times.\\
\midrule
26 & Measure an object with an instrument & The user agent must take a specific item (from 4 possible items), measure a specific property (from 5 possible instruments), and place the item in one of 4 containers depending on the range (e.g. temperature range 0-2.5C in Container A, temperature range 2.5-5.0C in Container B, etc.). \\
\midrule 
\rowcolor[HTML]{F3F3F3} 
27 & Pick-and-place object & The user agent must take a specific item (from 5 possible items) and place it in a specific container. \\
\midrule
28 & Pick-and-give object & The user must take a specific item (from 5 possible items) and give it to a specific NPC agent (from 3 possible agents). \\
\midrule
\rowcolor[HTML]{F3F3F3} 
29 & Read DiscoveryFeed posts & The user agent must take a specific item (from 5 possible items), and place it in a specific container.  The item to take is provided in a DiscoveryFeed post. \\
\midrule
30 & Move through doors & The user agent must successfully navigate a randomly generated maze (in the form of a house with walls/doors in random locations), and pick up a flag at the end of the maze.  All doors must be opened and the flag must be taken for task success. \\
\midrule
\rowcolor[HTML]{F3F3F3} 
31 & Using keys with doors & The user agent must successfully navigate a randomly generated house that contains 3 doors, and pick up a flag in the final room.  The doors require keys, which are randomly placed near the doors. \\
\midrule
32 & Navigate to a specific room in a house & The user agent must succesfully navigate to a specific room in a house (such as the \textit{kitchen} or \textit{bedroom}), which it must infer based on the contents of each room.  Once in the correct room, it signifies its selection by dropping a flag. \\
\midrule
\rowcolor[HTML]{F3F3F3} 
33 & Search an environment for an object & The user agent must navigate the interior and exterior of a house environment in search of a red flag. \\
\midrule
34 & Interact with a moving agent & The user agent must interact with 3 moving NPC agents spread throughout the environment, notifying each that it's about to rain, and they may want to return home to avoid becoming wet. \\
\bottomrule
\end{tabular}
\end{table}


\subsubsection{Unit Test Enviroments}
The unit tests are shown in Figure~\ref{fig:unit_test_screenshots}, and described in Table~\ref{tab:unit-task-descriptions}.




\subsection{Object Properties}
\label{sec:appendix-object-properties}

A list of common object properties provided by the base \textsc{Object} storage class is shown in Table~\ref{tab:object-properties}.  Note that specific objects (like a \textsc{QuantumCrystal}) may implement properties not included in the base class (such as \textit{resonanceFrequency}).

\begin{table}[t!]
\centering
\scriptsize
\caption{A list of object properties common to the \textsc{Object} storage class in \env. From this list is omitted any properties tied to a specific object, e.g., the fuel quantity in the rocket.}
\label{tab:object-properties}
\begin{tabular}{cp{0.20\linewidth}p{0.65\linewidth}}
\toprule
\textbf{\#} & \textbf{Property Name} & \textbf{Description} \\
\midrule
1 & grid location & Initial world location (X, Y) \\
2 & isMovable & Can it be moved? \\
3 & isPassable & Can an agent walk over this? \\
4 & obscuresObjectsBelow & Does it obscure/hide objects on layers below it? \\
5 & isActivatable & Is this a device? (more specifically, can it be activated/deactivated?) \\
6 & isActivated & Is this device currently activated? \\
7 & isUsable & Can this device be used with another object? (e.g., specifically through the 'use' action). \\
8 & isDialogable & Can it be dialoged with? \\
9 & isShovelable & Can it be shoveled? \\
10 & isReadable & Can it be read? \\
11 & document & Any text to read \\
12 & temperatureC & The default object temperature, in Celsius \\
13 & heatSourceMaxTemp & If it is a heat source, then this is the maximum temperature that it can reach \\
14 & coolSourceMinTemp & If it is a cool source, then this is the minimum temperature that it can reach \\
15 & isLiving & Is it alive? \\
16 & substanceName & Name of the substance \\
17 & isSubstance & Is it a substance? \\
18 & isAutoReacting & Does it react automatically with other substances? \\
19 & mixtureDict & Dictionary of substances and their proportions in the mixture \\
20 & requiresKey & If it requires a key to open/use, then this is a special ID for the key. If the value is <=0, then it doesn't require a key. \\
21 & keyID & If this object acts as a key, here's its ID (0 by default) \\
22 & radiocarbonAge & Radiocarbon dating age (in years). -1 means it's not applicable/inconclusive. \\
23 & radioisotopeValues & Radioisotope values. If empty, then it's not applicable/inconclusive. \\
24 & soilNutrients & Soil nutrients. If empty, then it's not applicable/inconclusive. \\
25 & needsNutrientLevels & For seeds/plants: What nutrient levels do they need to grow? \\
26 & antirequirementsNutrientLevels & A list of dictionaries, each containing a list of nutrient levels under which the seed/plant will NOT grow \\
27 & density & Object density (in g/cm³). <=0 means it's not applicable/inconclusive. \\
28 & microscopeModifierText & Modifier text, if the object is viewed under a microscope. This is a list of strings, which are displayed in the microscope view. \\
29 & microscopeDesc & Short description to be displayed under a microscope. \\
30 & color & Color description of the object. \\
31 & spectrum & Spectrum data of the object. \\
32 & ph & pH value of the object. \\
33 & radiationusvh & Radiation level in microsieverts per hour. \\
34 & nitrogen & Nitrogen content. \\
35 & phosphorus & Phosphorus content. \\
36 & potassium & Potassium content. \\
37 & cosCanBeLiquid & Can it exist in liquid form? \\
38 & cosCanBeSolid & Can it exist in solid form? \\
39 & cosCanBeGas & Can it exist in gas form? \\
40 & cosMeltingPointC & Melting point in Celsius. \\
41 & cosBoilingPointC & Boiling point in Celsius. \\
42 & cosCombustionPointC & Combustion point in Celsius. \\
43 & livingMinTemp & Minimum temperature for living organisms. \\
44 & livingMaxTemp & Maximum temperature for living organisms. \\
45 & isContainer & Is it a container? \\
46 & isOpenable & If it's a container, can you open/close it? \\
47 & isOpenContainer & If it's a container, then is it open? \\
48 & containerPrefix & Container prefix (e.g., "in" or "on"). \\
49 & isOpen & Closed by default. \\
50 & contentsVisible2D & If it is a container, do we render the contents in the 2D representation, or is that already handled (e.g., for pots/jars, that render generic contents if they contain any objects). \\
51 & contents & Contents of the container (other objects). \\
52 & parentContainer & Back-reference for the container that this object is in. \\
53 & parts & List of parts that this object is made of. \\
54 & subparts & Subparts of the object. \\
55 & isEdible & Can it be eaten? \\
56 & isCooked & Is it cooked? \\
57 & isPoisonous & Is it poisonous? \\
58 & isPassage & Is this a passage? \\
59 & materials & List of materials that this object is made of. \\
60 & manualMaterialNames & A list of material types to add during initialization (in code, rather than from the spreadsheet). \\
61 & actionHistory & Agent action history, if applicable (None by default). \\
62 & isAgent & Is this object an agent? (e.g., a person). \\
63 & isNPC & Is this agent an NPC? \\
\bottomrule
\end{tabular}
\end{table}



\clearpage
\section{Additional Baseline Model Details}
\label{sec:appendix-additional-model-details}

All models and their implementations are provided in the \env code repository. 

\subsection{ReAct}
The ReAct agent generates a thought and action at each step given a prompt describing the environment and the previous steps. The prompt contains information about the environment, how to interact with the environment, current environment state, teleportable locations, interactable objects, and valid actions. The previous steps are shown as a sequence of (thought, action, observation)s, e.g.:

\rule{\textwidth}{0.5pt}
\begin{footnotesize}
\begin{verbatim}
History of action-observations:
  Action:
    ```json
    {
    "action": "ROTATE_DIRECTION",
    "arg1": "north",
    "thought": "I need to face north to approach the door and enter the building."
    }
    ```
  Observation:
    ```json
    {"message": "I rotated to face north.", "errors": [], "success": true}
     ```
  Action:
    ```json
    {
    "action": "OPEN",
    "arg1": 2056,
    "thought": "I need to open the door in front of me to enter the building."
    }
    ```
  Observation:
    ```json
    {"message": "I opened the door.", "errors": [], "success": true}
    ```
\end{verbatim}
\end{footnotesize}
\rule{\textwidth}{0.5pt}

The agent generates thoughts and actions using the JSON formatted output as shown above. Each action is executed in the environment and the returned observation is then added to the trajectory. Since tasks in DiscoveryWorld can be extremely long, the trajectories can often exceed the maximum token limit of our LLM. We generally hit the 8K token limit of GPT4-O in 40 steps. To avoid this, we limit the trajectory in the prompt to a maximum of 10K characters by removing the oldest (thought, action, observation)s and replacing it with the string \texttt{[TRIMMED HISTORY]}.




\subsection{Plan-and-Execute}
Another approach to deal with long trajectories used in prior work~\cite{intercode,adapt} is to decompose the problem into a plan of simpler steps and executing each step independently using ReAct. Since it would be challenging to generate the entire plan in one shot \emph{without exploring the environment}, we use iterative decomposition~\cite{decomp} where we generate one step of the plan, then execute it and based on the result generate the next step. The planner uses the same base LLM with the prompt describing the format of the plan, e.g., 

\rule{\textwidth}{0.5pt}
\begin{footnotesize}
\begin{verbatim}
Task: Open the door and exit the building.

Plan:
Step 1: Go to the door that leads to the exit
 -- Task completed! I am at the door that leads to the exit
Step 2: Open the door
 -- Task failed! Door is locked and I don't have the key
Step 3: Find the key to the door
...
\end{verbatim}
\end{footnotesize}
\rule{\textwidth}{0.5pt}

To execute each step of the plan, we use the same ReAct agent as above but additionally add the previous steps of the plans and their results to the prompt. To prevent the ReAct agent from wasting time on infeasible steps, we add a hint to the prompt to consider returning the \texttt{Task failed!} message (e.g. Step 2 above) when $1/5^{th}$ of the environment step budget is used up. We use the same truncation strategy for the ReAct agent and no truncation is needed for the planner due to the much shorter plans. 


\subsection{Hypothesizer}
\label{sec:appendix-hypothesizer}
This agent maintains a working memory of science-related knowledge (allowing the knowledge this agent has discovered to be more directly evaluated).  To assist in planning and execution, this agent also maintains a running hypothesis for the task solution, as well as a short natural language plan for the immediate steps it should take.  More specifically, at each step, \textsc{Hypothesizer} chooses an action based on it's current working memory, plan, and running hypothesis.  It then reflects on the results of that action, and updates the working memory based on the results.  The working memory is science themed, storing two types of records: (1) \textsc{Hypotheses}, which include the \textit{hypothesis statement}, it's current \textit{status} (i.e. whether it has been confirmed, rejected, or is still pending), and a list of \textit{supporting evidence} the agent has listed as helping make this determination, (2) and \textsc{Measurements}, which take the form of specific observations the agent has made of objects in the environment -- for example, that a particular plant appeared to grow in a mixture including a specific nutrient.  This memory and reflection is similar to \textsc{CLIN}~\cite{Majumder2023CLINAC}, except: (a) this memory is structured to hold science-domain relations, where \textsc{CLIN}'s holds causal relations, and (b) this memory is updated after each action in the environment, where \textsc{CLIN}'s is updated only after running an entire task to completion.  This means that, if an agent takes $N$ steps in an environment, it requires $2N$ LLM calls to accommodate this \textit{act-then-reflect} cycle.  Because the memory can grow large, after every $10$ actions, the working memory is summarized through a separate LLM call that requests a maximum of $40$ separate \textit{hypothesis} or \textit{measurement} entries, and a maximum of $2k$ tokens for the entire memory.  When producing actions, \textsc{Hypothesizer} is prompted to produce (1) a specific action, (2) an explanation that includes its running hypothesis, and how this action helps move it towards evaluating this hypothesis, and (3) a short memory describing its plans or any other helpful information it wishes to pass on to the next step.  When prompted to produce actions, \textsc{Hypothesizer} is provided with (a) the current observation from the environment, (b) the last 3 actions it took (and their respective explanations and notes), (c) it's current working memory, (d) a list of possible actions and which arguments they take.  All output is produced as a \textsc{JSON} dictionary. 

The working memory that Hypothesizer stores can be long, especially for longer tasks. An example of the Hypothesizer agent's working memory at the successful completion of a specific shorter task (\textit{Proteomics, Normal, Seed 2}) is provided below.  

\begin{scriptsize}
\rule{\textwidth}{0.5pt}
\begin{verbatim}
Task Description: You are in a biological preserve on Planet X, that has 5 different animal species. 
We suspect that one of these animal species is not native to the area, but migrated from an isolated 
island in the recent past. Your task is to use the proteomics meter to analyze the proteins of each 
of the 5 animal species, which can be found throughout the environment away from the central statue 
area, and determine which species is the anomoly. Once you have completed your task, return to the 
statue area and drop the red flag directly beside the statue of the animal species that is the 
anomoly.

"working_memory": {
    "scientific_knowledge": [
        {
            "measurement": "The spheroid (uuid: 23195) is located in Possible Animal Area 1.",
            "step": 2
        },
        {
            "measurement": "The proteomics meter (uuid: 33276) is now in my inventory.",
            "step": 8
        },
        {
            "measurement": "The proteomics meter (uuid: 33276) requires two argument objects to be 
               used.",
            "step": 10
        },
        {
            "hypothesis": "if (object:spheroid) is (action:directly in front) then 
               (object:proteomics meter) can be used to analyze its proteins.",
            "status": "confirmed",
            "step": 13,
            "supporting evidence": "Using the proteomics meter on the spheroid yielded protein data: 
               Protein A: 0.87, Protein B: 0.80."
        },
        {
            "measurement": "The spheroid (uuid: 23195) has Protein A level of 0.87 and Protein B 
               level of 0.80 when analyzed using the proteomics meter.",
            "step": 13
        },
        {
            "measurement": "The animaplant (uuid: 20030) is located in Possible Animal Area 2.",
            "step": 14
        },
        {
            "hypothesis": "if (object:animaplant) is (action:analyzed) using (object:proteomics 
               meter) then (property:protein data) will be obtained.",
            "status": "confirmed",
            "step": 15,
            "supporting evidence": "Using the proteomics meter on the animaplant yielded protein 
               data: Protein A: 0.76, Protein B: 0.67."
        },
        {
            "measurement": "The animaplant (uuid: 20030) has Protein A level of 0.76 and Protein B 
               level of 0.67 when analyzed using the proteomics meter.",
            "step": 15
        },
        {
            "measurement": "The animal in Possible Animal Area 3 (uuid: 20031) is not accessible.",
            "step": 17
        },
        {
            "hypothesis": "if (object:animal) is (action:teleported to) using (object:uuid: 20031) 
               then (property:accessibility) will be false.",
            "status": "confirmed",
            "step": 22,
            "supporting evidence": "Attempting to teleport to the animal object with UUID 20031 
               resulted in an error indicating that no object with that UUID was found."
        },
        {
            "measurement": "There is a tree (big) trunk blocking the path to the west.",
            "step": 25
        },
        {
            "measurement": "I successfully moved south to navigate around the tree (big) trunk.",
            "step": 26
        },
        {
            "measurement": "I successfully moved west to get directly in front of the pink 
               animal.",
            "step": 28
        },
        {
            "measurement": "I rotated to face west, aligning myself with the pink animal.",
            "step": 30
        },
        {
            "measurement": "The object directly in front of me, which I thought was a pink 
               animal, is actually grass (uuid: 20981).",
            "step": 31
        },
        {
            "measurement": "The proteomics meter (uuid: 33276) was used to investigate the grass 
               (uuid: 20981), but the results were inconclusive.",
            "step": 31
        },
        {
            "measurement": "The object with UUID 44261 is a plant, not an animal.",
            "step": 37
        },
        {
            "measurement": "The proteomics meter (uuid: 33276) was used to investigate the 
               animaplant (uuid: 44261). The results are as follows: 
               Protein A: 0.76, Protein B: 0.67.",
            "step": 37
        },
        {
            "measurement": "The object with UUID 12545 is a plant, not an animal.",
            "step": 38
        },
        {
            "measurement": "The proteomics meter (uuid: 33276) was used to investigate the plant
               (uuid: 12545), but the results were inconclusive.",
            "step": 38
        },
        {
            "measurement": "The vortisquid (uuid: 41429) is one of the animal species in the 
               biological preserve.",
            "step": 149
        },
        {
            "hypothesis": "if (object:vortisquid) is (action:analyzed) using 
               (object:proteomics meter) then (property:protein data) will be obtained.",
            "status": "confirmed",
            "step": 150,
            "supporting evidence": "Using the proteomics meter on the vortisquid yielded 
               protein data: Protein A: 0.38, Protein B: 0.46."
        },
        {
            "measurement": "The vortisquid (uuid: 41429) has Protein A level of 0.38 and 
               Protein B level of 0.46 when analyzed using the proteomics meter.",
            "step": 150
        },
        {
            "measurement": "The prismatic beast (uuid: 25282) is located in Possible Animal 
               Area 5.",
            "step": 151
        },
        {
            "hypothesis": "if (object:prismatic beast) is (action:analyzed) using 
               (object:proteomics meter) then (property:protein data) will be obtained.",
            "status": "confirmed",
            "step": 153,
            "supporting evidence": "Using the proteomics meter on the prismatic beast 
               yielded protein data: Protein A: 0.73, Protein B: 0.68."
        },
        {
            "measurement": "The prismatic beast (uuid: 25282) has Protein A level of 0.73 
               and Protein B level of 0.68 when analyzed using the proteomics meter.",
            "step": 153
        },
        {
            "measurement": "The prismatic beast (uuid: 5059) is located in Possible Animal 
               Area 7.",
            "step": 156
        },
        {
            "measurement": "The prismatic beast (uuid: 5059) has Protein A level of 0.73 and 
               Protein B level of 0.68 when analyzed using the proteomics meter.",
            "step": 157
        },
        {
            "measurement": "The echojelly (uuid: 4429) is one of the animal species in the 
               biological preserve.",
            "step": 160
        },
        {
            "hypothesis": "if (object:echojelly) is (action:analyzed) using 
               (object:proteomics meter) then (property:protein data) will be obtained.",
            "status": "confirmed",
            "step": 161,
            "supporting evidence": "Using the proteomics meter on the echojelly yielded 
               protein data: Protein A: 0.86, Protein B: 0.82."
        },
        {
            "measurement": "The echojelly (uuid: 4429) has Protein A level of 0.86 and 
               Protein B level of 0.82 when analyzed using the proteomics meter.",
            "step": 161
        },
        {
            "measurement": "The vortisquid (uuid: 41429) has significantly different protein 
               levels (Protein A: 0.38, Protein B: 0.46) compared to the other species, 
               indicating it is the anomaly.",
            "step": 162
        },
        {
            "measurement": "The object with UUID 33979 (red flag) is not directly in front of 
               me or in my inventory.",
            "step": 172
        },
        {
            "hypothesis": "if (object:red flag) is (action:teleported to) using 
               (object:uuid: 33979) then (property:accessibility) will be true.",
            "status": "pending",
            "step": 172,
            "supporting evidence": "Attempting to pick up the red flag from the Instrument 
               Table resulted in errors indicating that the object with UUID 33979 could 
               not be found."
        },
        {
            "measurement": "The object with UUID 33979 (red flag) is not directly in front 
               of me or in my inventory.",
            "step": 173
        },
        {
            "hypothesis": "if (object:red flag) is (action:teleported to) using 
               (object:uuid: 33979) then (property:accessibility) will be true.",
            "status": "pending",
            "step": 173,
            "supporting evidence": "Attempting to pick up the red flag from the Instrument 
               Table resulted in errors indicating that the object with UUID 33979 could 
               not be found."
        },
        {
            "measurement": "The object with UUID 33979 (red flag) is located at 
               coordinates (15, 20).",
            "step": 174
        },
        {
            "hypothesis": "if (object:red flag) is (action:teleported to) using 
               (object:uuid: 33979) then (property:accessibility) will be true.",
            "status": "confirmed",
            "step": 174,
            "supporting evidence": "Successfully teleported to the object with 
               UUID 33979 (red flag) at coordinates (15, 20)."
        },
        {
            "measurement": "The object with UUID 33979 (red flag) is now in my inventory.",
            "step": 175
        },
        {
            "hypothesis": "if (object:red flag) is (action:picked up) then 
               (property:accessibility) will be true.",
            "status": "confirmed",
            "step": 175,
            "supporting evidence": "Successfully picked up the red flag after teleporting 
               to it."
        },
        {
            "measurement": "The object with UUID 33979 (red flag) is located at 
               coordinates (15, 20).",
            "step": 176
        },
        {
            "hypothesis": "if (object:red flag) is (action:teleported to) using 
               (object:uuid: 33979) then (property:accessibility) will be true.",
            "status": "confirmed",
            "step": 176,
            "supporting evidence": "Successfully teleported to the object with 
               UUID 33979 (red flag) at coordinates (15, 20)."
        },
        {
            "measurement": "The object with UUID 33979 (red flag) is now in my inventory.",
            "step": 176
        },
        {
            "hypothesis": "if (object:red flag) is (action:picked up) then 
               (property:accessibility) will be true.",
            "status": "confirmed",
            "step": 176,
            "supporting evidence": "Successfully picked up the red flag after teleporting to it."
        },
        {
            "measurement": "I teleported to the Statue of a vortisquid at 
               coordinates (18, 17).",
            "step": 176
        }
    ]
}    
\end{verbatim}
\end{scriptsize}
\rule{\textwidth}{0.5pt}

\subsection{Cost analysis}
\label{sec:appendix_cost_analysis}

An approximate cost analysis of each baselines model is provided below, in Table~\ref{tab:cost-estimate}.  Model cost varies across task, environment complexity, and the number of items stored within a model's memory (e.g. the size of the history in \textsc{ReAct}, or the size of the memory in \textsc{Hypothesizer}), so approximate averages are shown. 

The full benchmark contains $8$ discovery themes $\times$ $3$ difficulty levels $\times$ $5$ seeds $= 120$ tasks to evaluate.  \textit{Easy} tasks are evaluated to 100 steps, while \textit{Normal} and \textit{Challenge} tasks are evaluated to 1000 steps or a hard \$125 limit, whichever came first.  This results in a total upper-bound estimate of $40\times100 (easy) + 40\times1000 (normal) + 40\times1000(challenge) = 84,000$ steps required to complete an evaluation run, assuming no agents finish early due to task completion.

\begin{table}[h!]
    \centering
    \footnotesize
    \caption{Approximate cost estimate for the models investigated in this work (all using GPT-4o). Assumes a cost of \$5/M input tokens, and \$15/M output tokens (current pricing as of this writing).}
    \label{tab:cost-estimate}
    \begin{tabular}{ccc}
    \toprule
    \textbf{Model}  &   \textbf{Cost per 100 steps}  &   \textbf{Total cost estimate (120 tasks, 84k steps)} \\
    \midrule
        \textsc{ReAct}               &  \$2-4 / 100 steps       &   \$3,360 \\
        \textsc{Plan + Execute}      &  \$3-4 / 100 steps       &   \$3,360 \\
        \textsc{Hypothesizer}        &  \$10  / 100 steps       &   \$8,400 \\
    \bottomrule
    \end{tabular}
\end{table}

\section{Automatic Evaluation of Explanatory Discovery Knowledge}
\label{sec:appendix-automatic-knowledge-evaluation}

The following prompt is used to automatically evaluate explanatory discovery knowledge, in reference to gold reference questions provided automatically by the \env task scorer:

\rule{\textwidth}{0.5pt}
\begin{footnotesize}
\begin{verbatim}
Your task is to evaluate the knowledge of an agent playing a game about 
scientific discovery, to see if they have figured out the discovery 
required to solve the task. I will supply you with: (1) the task 
description, (2) the agent's knowledge, and (3) a single critical question
about the discovery, that you should answer with either true (1) or 
false (0). 

Task Description:
```
[*INSERT TASK DESCRIPTION HERE*]
```

Agent's Knowledge:
```
[*INSERT AGENT'S KNOWLEDGE HERE*]
```

Critical Question:
```
[*INSERT SINGLE KNOWLEDGE EVALUATION QUESTION HERE FROM SCORECARD*]
```

Please answer this question by responding `1` if the agent's knowledege 
reflects having discovered the information in the critical question, and
`0` otherwise.  This response should be in the `evaluation` key of the
response. The response format is a JSON dictionary containing three keys:
`criticalQuestion`, `evaluation`, and `explanation`. 
```
{
    \"criticalQuestion\": \"repeat the critical question\",
    \"evaluation\": 0 or 1 (as integers),
    \"explanation\": \"provide a brief explanation for evaluation, making 
       reference to the agent's knowledge and whether or not it reflects 
       the critical question.\"
}
```
\end{verbatim}
\end{footnotesize}
\rule{\textwidth}{0.5pt}

The following is an example of the scorer producing a \textbf{positive assessment}: 

\rule{\textwidth}{0.5pt}
\begin{footnotesize}
\begin{verbatim}
Task: Proteomics

"evaluation": [
    {
        "criticalQuestion": "Does it clearly state that the vortisquid has 
              proteomics values that appear to be outliers compared to the 
              other animals, whose proteomics values all cluster together?",
        "evaluation": 1,
        "explanation": "The agent's knowledge includes a specific 
              measurement stating that the vortisquid has significantly 
              different protein levels (Protein A: 0.38, Protein B: 0.46) 
              compared to the other species, indicating it is the anomaly. 
              This directly addresses the critical question by identifying 
              the vortisquid's proteomics values as outliers compared to 
              the other animals."
    }
],
"evaluation_totalscore_raw": 1,
"evaluation_totalscore": 1.0,
\end{verbatim}
\end{footnotesize}
\rule{\textwidth}{0.5pt}

The following is an example of the scorer producing a \textbf{negative assessment}:

\rule{\textwidth}{0.5pt}
\begin{footnotesize}
\begin{verbatim}
Task: Translation 

"evaluation": [
    {
        "criticalQuestion": "Does it clearly state that: The word 'Womple'
              means something like 'bring me'?",
        "evaluation": 0,
        "explanation": "The agent's knowledge does not explicitly state 
              that the word 'Womple' means 'bring me'. The agent has 
              recorded instances where the elder and the trainer used the 
              word 'Womple', but there is no clear interpretation or 
              translation of the word provided in the knowledge. The agent
              has noted that the elder seems interested in items in the 
              inventory and asked if the agent wants to give the items to 
              him, but this does not directly translate 'Womple' to 
              'bring me'."
    },
    {
        "criticalQuestion": "Does it clearly state that: The word 
              'florpt' means 'flower'?",
        "evaluation": 0,
        "explanation": "The agent's knowledge does not explicitly state 
              that the word 'florpt' means 'flower'. While there are 
              measurements and interactions involving the word 'florpt' 
              and objects like flowers, there is no direct confirmation 
              or hypothesis that equates 'florpt' to 'flower'."
    }
],
"evaluation_totalscore_raw": 0,
"evaluation_totalscore": 0.0,
\end{verbatim}
\end{footnotesize}
\rule{\textwidth}{0.5pt}

\section{Human Scientist Participants}
\label{sec:appendix_human}

\paragraph{General Details} To compare model performance against human performance, we recruited 11 practicing human scientists to complete the \env tasks, with their performance shown in Table~\ref{tab:human-performance}.  Scientists were recruited on the \textsc{UpWork} platform, each with: (1) an \textsc{MSc} or \textsc{PhD} in a natural science, (2) self-evaluated comfort and fluency with statistical methods and common software like spreadsheets, (3) comfort and previous experience with 2D top-down games.  To evaluate the latter, all participants were required to complete a screening task involving a \textit{Tutorial} in \env. Participants were given a maximum of one hour to complete each task.  While our agent models are run in a zero-shot setting, humans can't forget their past experiences, so our participants only completed the \textit{normal} and \textit{challenge} settings, to prevent any easily-discovered knowledge in the \textit{easy} setting from providing an advantage in other settings.  To collect accurate estimates of task difficulty, all humans completed the same seed (\textsc{Seed 0}) of each task.  To evaluate discovered knowledge, participants were asked to keep notes, and explicitly state their hypotheses, and supporting evidence.  Due to the varied nature of these notes (some wrote a few words, others entire pages per task) and their modality (text, spreadsheets), discovery knowledge was evaluated manually using the following criteria: $1$ (discovered all discovery knowledge in scorecard), $0.5$ (discovered some discovery knowledge in scorecard), or $0$ (did not find critical discovery knowledge).  Other metrics \textit{(task completion, procedural knowledge)} were evaluated automatically.  At their option, not all participants completed all tasks. 

\paragraph{Informed Consent} Details about the project's objective and the intended use of data generated during game play was provided to participants in the initial job description. Participants were further provided with a participation agreement detailing that all data generated from the project will be owned by AI2, which they agreed to by choosing to participate in the project.

\paragraph{Participant Instructions} Participant instructions can be found in  \href{https://github.com/allenai/discoveryworld/blob/main/README-USERSTUDY.md}{README-USERSTUDY.MD} on the \env code repository.

\paragraph{Potential Participant Risks} To protect the personal information and anonymity of the human scientists we recruited, no personally identifying information was collected and all data is anonymized.

\paragraph{Recruitment} Participants were recruited and hired on Upwork, where they were initially required to submit a proposal detailing their educational background, experience with statistical analyses and Python programming, and interest in video games. 

\paragraph{Compensation} Human scientists were compensated at a rate ranging from USD\$20-\$30/hr for their work, based on the rate they bid in their initial job proposals. The total amount spent on participant compensation did not exceed \$4,000.

\paragraph{Participants' Expertise} The participants we recruited had Master’s or Doctorate degrees in a range of disciplines. Their educational and/or work experience spanned the following fields: astrophysics, bioinformatics, biology, biomedical engineering, biostatistics, chemistry, computer engineering, data science, electrical engineering, geology, machine learning, mathematics, and physics.

\paragraph{Graphical Interface} Humans made use of the 2D graphical user interface for their study, rather than a text-only version (as some agent models use). Figure~\ref{fig:GUI} shows an example of the interface.

\paragraph{Time limit} Human participants were given a maximum of 1 hour to complete a given task. 

\paragraph{Zero-shot Performance} While participants were allowed to retry tasks that they did not complete successfully the first time, we did not include any retries in our evaluation of human performance, to give an accurate analog of \textit{zero-shot} model performance.

\paragraph{Data} Instructions on how to get the anonymized data from human participants is available at \href{https://github.com/allenai/discoveryworld}{https://github.com/allenai/discoveryworld}. 

\paragraph{Save failures} \env automatically saves log files to evaluate performance.  Some participants noted infrequent issues with this automatic save on the \textsc{Windows} platform that we were unable to resolve. As such, a small amount of data was lost, and is not included in this analysis. 

\begin{figure}
    \centering
    \includegraphics[width=0.75\textwidth]{figures/frame_1.png}
    \caption{Example of the user interface the participants used. This is for the \textit{It's (not) Rocket Science!} theme.}
    \label{fig:GUI}
\end{figure}









